\chapter{Hukum Bilangan Besar dan Teorema Limit Pusat}

Satu program pendek menghitung rata-rata dari $\exp(U)$ untuk sepuluh ribu bilangan
acak $U$ yang seragam pada $(0,1)$. Hasil cetaknya $1{,}722206$. Nilai sebenarnya
adalah $e-1=1{,}7182818285$, jadi galat jalankan itu $+0{,}003924$. Program
dijalankan ulang dengan empat kali lipat bilangan: $1{,}717502$, galat
$-0{,}000780$. Dengan sepuluh kali lipat: $1{,}717767$, galat $-0{,}000515$.
Polanya tidak rapi. Empat kali data memangkas galat lebih dari lima kali, sedangkan
sepuluh kali data hanya memangkasnya $7{,}6$ kali. Tebakan naif "sepuluh kali data,
sepersepuluh galat" tidak berlaku pada satu jalankan. Satu jalankan program tidak
memuat laju apa pun; yang punya laju adalah sebaran dari banyak jalankan.

Sebaran itu bisa diukur. Dengan $n=10^4$ simpangan baku penaksir ini $0{,}00492$;
dengan $n=4\cdot10^4$ menjadi $0{,}00246$; dengan $n=10^5$ menjadi $0{,}00156$.
Setiap kali $n$ naik empat kali, simpangan itu tepat separuh. Di sini ada dua
pernyataan berbeda yang sering tertukar. Pertama, rata-rata memang mendekat ke
nilai sebenarnya: itu kecenderungan, dan namanya hukum bilangan besar. Kedua, sisa
galat punya ukuran dan bentuk tertentu: itu fluktuasi, dan namanya teorema limit
pusat. Bab ini memberi nama, bukti, dan batas keberlakuan untuk keduanya.

Bab 18 sudah menyiapkan bahan: harapan, variansi, dan distribusi jumlah lewat
konvolusi. Konvolusi hanya berguna bila $n$ kecil dan bentuk asalnya sederhana;
untuk sebagian besar distribusi tidak ada bentuk tertutup. Yang kita butuhkan
bukan distribusi tepat pada satu $n$, melainkan pernyataan yang berlaku untuk
semua $n$ yang besar. Pertanyaan praktisnya dua: seberapa dekat rata-rata boleh
dipercaya, dan seberapa besar galat yang wajar. Jawabannya datang dari
memeriksa bagaimana barisan distribusi itu berperilaku, bukan dari menghitung
satu distribusi.

\begin{untukanak}
Bayangkan satu kelas diminta menebak banyaknya kelereng di dalam stoples besar.
Tebakan anak-anak jauh meleset ke dua arah. Ada yang menebak tiga ratus, ada yang
menebak dua ribu, padahal isinya delapan ratus. Guru lalu menuliskan semua tebakan
di papan dan menghitung rata-ratanya. Angkanya ternyata dekat sekali dengan
delapan ratus.

Sekarang gurunya menambah murid. Dengan tiga puluh penebak, rata-ratanya masih
goyang. Dengan tiga ratus penebak, rata-ratanya lebih tenang. Dengan tiga ribu
penebak, rata-ratanya hampir tidak bergeser lagi. Ada dua hal aneh di sini.
Pertama, tidak seorang pun penebak yang menjadi lebih jago; yang berubah hanya
banyaknya penebak. Kedua, menambah penebak dari tiga puluh ke tiga ratus menolong
jauh lebih banyak daripada menambah dari tiga ratus ke tiga ribu. Pertolongannya
terus bertambah, tetapi makin lama makin sedikit.

Pertanyaan yang membuat penasaran: kalau setiap tebakan bisa meleset sangat jauh,
mengapa rata-ratanya justru bisa dipercaya? Dan mengapa pertolongannya melambat
dengan pola yang bisa ditebak, bukan sembarangan? Di bab ini kita akan tahu
mengapa, dan sekaligus tahu kapan kebiasaan mempercayai rata-rata itu menipu.
\end{untukanak}

\section{Tiga arti kata menuju pada variabel acak}
\tingkat{dasar}

Kalimat sehari-harinya: barisan variabel acak bisa "menuju" dalam tiga arti yang
berbeda, dan arti yang paling longgar sama sekali tidak mengatakan bahwa nilainya
mendekat.

Barisan bilangan real hanya punya satu arti menuju, yaitu jarak $|x_n-L|$ menjadi
kecil. Untuk variabel acak, jarak itu sendiri adalah variabel acak, sehingga kita
harus memilih besaran mana yang dibuat kecil. Pilihan pertama adalah peluangnya.
Ini pilihan yang paling lemah secara intuitif tetapi paling berguna secara praktis.

\kamus{konvergensi dalam peluang}{convergence in probability}{janji bahwa peluang
meleset lebih dari toleransi mengecil menuju nol, bukan janji bahwa nilainya
mendekat pada satu jalankan}

\begin{definisi}[Konvergensi dalam peluang]
Barisan variabel acak $X_1,X_2,\dots$ \emph{konvergen dalam peluang}
\index{konvergensi dalam peluang} ke variabel acak $X$ jika untuk setiap
$\varepsilon>0$ berlaku
\[
P\bigl(|X_n-X|>\varepsilon\bigr)\to 0 \qquad (n\to\infty).
\]
\end{definisi}

Perhatikan letak kuantornya. Nilai $\varepsilon$ dipilih lebih dahulu, lalu
peluangnya boleh mengecil sesudah itu. Definisi ini tidak mengatakan bahwa
$|X_n-X|$ mengecil untuk satu hasil percobaan tertentu; ia mengatakan bahwa
kejadian "meleset lebih dari $\varepsilon$" menjadi makin jarang. Pada percobaan
yang sama, $X_n$ masih boleh meleset jauh, asal kejadian itu jarang.

\kamus{konvergensi dalam distribusi}{convergence in distribution}{yang menyerupai
adalah bentuk sebarannya, bukan nilainya; dua variabel bisa sebentuk tanpa pernah
berdekatan}

\begin{definisi}[Konvergensi dalam distribusi]
Misalkan $F_n$ fungsi sebaran $X_n$ dan $F$ fungsi sebaran $X$. Kita tulis
$X_n\rightsquigarrow X$ bila $F_n(t)\to F(t)$ untuk setiap titik $t$ tempat $F$
kontinu. Batasnya disebut \emph{limit dalam distribusi}\index{konvergensi dalam
distribusi}.
\end{definisi}

Syarat "di titik kontinu" bukan keisengan teknis. Ambil $X_n$ yang bernilai $1/n$
dengan peluang satu, dan $X$ yang bernilai $0$ dengan peluang satu. Maka $F_n(0)=0$
untuk setiap $n$, sedangkan $F(0)=1$, jadi barisan itu tidak konvergen di $t=0$.
Di setiap $t\neq0$ barisan itu konvergen, dan $t=0$ adalah satu-satunya titik
ketakkontinuan $F$. Tanpa syarat tadi, kesimpulan yang benar akan ditolak hanya
karena satu titik.

Arti ketiga lebih kuat dan lebih harfiah: hampir setiap hasil percobaan, sebagai
satu titik $\omega$ di ruang sampel, menghasilkan barisan bilangan real yang
konvergen.

\begin{definisi}[Konvergensi hampir pasti]
$X_n$ \emph{konvergen hampir pasti}\index{konvergensi hampir pasti} ke $X$ bila
\[
P\bigl(\{\omega : X_n(\omega)\to X(\omega)\}\bigr)=1 .
\]
\end{definisi}

Di sini barisan bilangan real $X_1(\omega),X_2(\omega),\dots$ benar-benar
konvergen dalam arti Bab 4, dan himpunan hasil percobaan yang gagal konvergen
hanya boleh berpeluang nol. Karena itu konvergensi hampir pasti berbicara tentang
satu jalankan yang panjang, sedangkan konvergensi dalam peluang berbicara tentang
satu $n$ pada banyak jalankan sekaligus.

Ketiga arti itu tidak setara, tetapi ada arah implikasi yang tetap.

\begin{teorema}[Hubungan tiga modus]
Untuk variabel acak $X_n$ dan $X$ berlaku:
(a) $X_n\to X$ hampir pasti mengakibatkan $X_n\to X$ dalam peluang;
(b) $X_n\to X$ dalam peluang mengakibatkan $X_n\rightsquigarrow X$;
(c) bila $X_n\rightsquigarrow c$ dan $c$ konstanta, maka $X_n\to c$ dalam peluang.
\end{teorema}

\begin{proof}[Bukti]
Untuk (a), tetapkan $\varepsilon>0$ dan misalkan
$A_n=\{\sup_{k\ge n}|X_k-X|>\varepsilon\}$. Himpunan $A_n$ menurun, dan setiap
$\omega$ di $\bigcap_n A_n$ memenuhi $|X_k(\omega)-X(\omega)|>\varepsilon$ untuk
tak hingga banyak $k$, sehingga barisan itu tidak konvergen ke $X(\omega)$.
Jadi $\bigcap_n A_n$ termuat di himpunan berpeluang nol. Karena peluang bersifat
kontinu dari atas pada himpunan berpeluang hingga,
$P(A_n)\to P(\bigcap_n A_n)=0$. Kejadian $\{|X_n-X|>\varepsilon\}$ termuat di
$A_n$, maka peluangnya juga menuju nol.

Untuk (b), ambil $t$ titik kontinu $F$ dan $\varepsilon>0$. Dari
$\{X_n\le t\}\subseteq\{X\le t+\varepsilon\}\cup\{|X_n-X|>\varepsilon\}$ kita
peroleh
\[
F_n(t)\le F(t+\varepsilon)+P\bigl(|X_n-X|>\varepsilon\bigr).
\]
Dengan cara yang sama, dari $\{X\le t-\varepsilon\}\subseteq\{X_n\le t\}\cup
\{|X_n-X|>\varepsilon\}$ kita peroleh
$F(t-\varepsilon)\le F_n(t)+P(|X_n-X|>\varepsilon)$. Kedua ketaksamaan itu
memberi
\[
F(t-\varepsilon)-\delta_n\le F_n(t)\le F(t+\varepsilon)+\delta_n,
\qquad \delta_n=P\bigl(|X_n-X|>\varepsilon\bigr).
\]
Karena $\delta_n\to0$, limit inferior dan limit superior $F_n(t)$ terkurung di
antara $F(t-\varepsilon)$ dan $F(t+\varepsilon)$. Biarkan $\varepsilon\downarrow0$
dan pakai kekontinuan $F$ di $t$; keduanya menuju $F(t)$.

Untuk (c), ambil $\varepsilon>0$. Karena $F$ adalah sebaran titik massa di $c$,
titik $c-\varepsilon$ dan $c+\varepsilon$ keduanya titik kontinu $F$, sehingga
\[
P\bigl(|X_n-c|>\varepsilon\bigr)
\le P(X_n<c-\varepsilon)+P(X_n>c+\varepsilon)
\le F_n(c-\varepsilon)+1-F_n(c+\varepsilon)
\to 0+1-1=0 .
\]
\end{proof}

Tiga arah sebaliknya tidak berlaku, dan contoh penyangkalnya perlu diingat karena
masing-masing menunjukkan apa yang sesungguhnya hilang.

\begin{contoh}[Dalam peluang tetapi tidak hampir pasti]
Pecah $(0,1)$ menjadi $k$ bagian sama panjang, untuk $k=1,2,3,\dots$, lalu
urutkan semua bagian itu menjadi satu daftar. Definisikan $X_n(\omega)=1$ bila
$\omega$ berada di bagian ke-$n$ dan $0$ bila tidak. Bagian ke-$n$ berasal dari
pemecahan dengan $k_n$ bagian, jadi $P(X_n=1)=1/k_n\to0$ dan $X_n\to0$ dalam
peluang. Tetapi untuk setiap $\omega\in(0,1)$, dalam setiap blok pemecahan ada
tepat satu $n$ dengan $X_n(\omega)=1$. Nilai satu itu muncul tak hingga sering,
sehingga $X_n(\omega)$ tidak konvergen, untuk setiap $\omega$. Jadi tidak ada
konvergensi hampir pasti. Barisan ini bahkan konvergen dalam $L^1$, karena
$\mathbb{E}|X_n|=1/k_n\to0$, maka kekuatan $L^1$ pun tidak menutup celah itu.
\end{contoh}

\begin{contoh}[Dalam distribusi tetapi tidak dalam peluang]
Ambil $X\sim N(0,1)$ dan tetapkan $X_n=-X$ untuk setiap $n$. Sebaran $X_n$ sama
persis dengan sebaran $X$, jadi $F_n(t)=F(t)$ dan $X_n\rightsquigarrow X$ secara
hampa. Tetapi $P(|X_n-X|>\varepsilon)=P(2|X|>\varepsilon)>0$ untuk setiap
$\varepsilon$, dan peluang itu tidak bergantung pada $n$. Dua variabel yang
sebentuk belum tentu berdekatan: $X$ dan $-X$ bisa berjarak jauh pada jalankan
yang sama.
\end{contoh}

Pelajaran dari dua contoh itu: konvergensi dalam distribusi hanya menyentuh
sebaran satu per satu, sedangkan dua modus lainnya menyentuh hubungan
antar-variabel pada ruang sampel yang sama. Karena itu teorema (c) penting secara
praktis: dalam statistika limit kita hampir selalu konstanta, dan di situ ketiga
modus bertemu.

Satu peringatan lain harus dicatat sejak awal, karena kesalahannya mahal.

\begin{catatan}
Konvergensi dalam peluang tidak membawa serta konvergensi harapan. Misalkan
$P(X_n=n^2)=1/n$ dan $P(X_n=0)=1-1/n$. Untuk setiap $\varepsilon>0$,
$P(|X_n|>\varepsilon)=1/n\to0$, jadi $X_n\to0$ dalam peluang. Tetapi
$\mathbb{E}[X_n]=n$, sehingga harapannya menuju tak hingga. Rata-rata yang
memusat bukan jaminan bahwa harapannya memusat.
\end{catatan}

\section{Ketaksamaan Chebyshev dan hukum bilangan besar lemah}
\tingkat{menengah}

Kalimat sehari-harinya: peluang menyimpang jauh dari rata-rata dibatasi oleh
variansi dibagi kuadrat jarak penyimpangan itu.

Untuk membuktikan hukum bilangan besar kita hanya butuh satu ketaksamaan, dan
ketaksamaan itu butuh satu pengamatan sederhana tentang harapan.

\begin{lemma}[Markov]
Jika $Y\ge0$ dan $a>0$, maka $P(Y\ge a)\le \mathbb{E}[Y]/a$.
\index{ketaksamaan Markov}
\end{lemma}

\begin{proof}[Bukti]
Tulis $Y=Y\mathbf{1}_{\{Y<a\}}+Y\mathbf{1}_{\{Y\ge a\}}$. Suku pertama tak
negatif. Pada kejadian $Y\ge a$ berlaku $Y\ge a$, sehingga harapan suku kedua
paling sedikit $a\,P(Y\ge a)$. Jadi
$\mathbb{E}[Y]\ge 0+a\,P(Y\ge a)$.
\end{proof}

\begin{akibat}[Ketaksamaan Chebyshev]
Jika $\mathbb{E}[X]=\mu$ dan $\operatorname{Var}(X)=\sigma^2<\infty$, maka untuk
setiap $a>0$
\[
P\bigl(|X-\mu|\ge a\bigr)\le\frac{\sigma^2}{a^2}.
\]
\index{ketaksamaan Chebyshev}
\end{akibat}

\begin{proof}[Bukti]
Terapkan Lemma Markov pada $Y=(X-\mu)^2$ dengan ambang $a^2$. Kejadian
$\{Y\ge a^2\}$ sama dengan $\{|X-\mu|\ge a\}$, dan
$\mathbb{E}[Y]=\sigma^2$.
\end{proof}

Ketaksamaan ini lemah, dan kelemahannya bukan cacat tersembunyi melainkan
kekuatannya. Ia tidak mengetahui bentuk distribusi apa pun, hanya dua momen.
Karena itu ia berlaku untuk semua distribusi bervariansi hingga, dan justru
karena itu ia cukup untuk membuktikan teorema yang berlaku umum.

\kamus{hukum bilangan besar}{law of large numbers}{janji bahwa rata-rata banyak
pengamatan menempel ke rata-rata sebenarnya, dengan peluang meleset yang mengecil}

\begin{teorema}[Hukum bilangan besar lemah]
Misalkan $X_1,X_2,\dots$ bebas dan berdistribusi sama dengan
$\mathbb{E}[X_i]=\mu$ dan $\operatorname{Var}(X_i)=\sigma^2<\infty$. Maka untuk
setiap $\varepsilon>0$\index{hukum bilangan besar}
\[
P\bigl(|\bar X_n-\mu|\ge\varepsilon\bigr)\to0,
\qquad \bar X_n=\frac1n\sum_{i=1}^n X_i .
\]
\end{teorema}

\begin{proof}[Bukti]
Karena harapan linear, $\mathbb{E}[\bar X_n]=\mu$. Karena variabelnya bebas,
variansi jumlah adalah jumlah variansi, sehingga
$\operatorname{Var}(\bar X_n)=\sigma^2/n$. Ketaksamaan Chebyshev dengan
$a=\varepsilon$ memberi
\[
P\bigl(|\bar X_n-\mu|\ge\varepsilon\bigr)\le\frac{\sigma^2}{n\varepsilon^2}.
\]
Ruas kanan menuju nol untuk setiap $\varepsilon>0$ tetap.
\end{proof}

Perhatikan apa yang baru saja dibuktikan dan apa yang tidak. Yang dibuktikan:
untuk toleransi $\varepsilon$ yang kita pilih sendiri, peluang rata-rata meleset
lebih dari $\varepsilon$ menuju nol. Yang tidak dibuktikan: seberapa cepat, dan
bagaimana bentuk sisa galatnya. Keduanya pertanyaan berbeda.

Gambar~\ref{fig:tiga-lebar} menggambar isi pernyataan itu untuk dadu. Yang
berubah dari $n=4$ ke $n=64$ hanya dua: lebarnya menyusut, puncaknya meninggi,
sedangkan pusatnya tetap di $3{,}5$. Hukum bilangan besar bukan janji bahwa satu
jalankan lebih tepat, melainkan bahwa sebaran kemungkinannya mengerut ke satu
titik.

\begin{figure}[htbp]
\centering
\begin{tikzpicture}
\begin{axis}[
  width=0.95\linewidth, height=5.4cm,
  xmin=1, xmax=6, ymin=0, ymax=2.05,
  xlabel={nilai rata-rata $\bar x_n$}, ylabel={kerapatan},
  label style={font=\small}, tick label style={font=\footnotesize},
  axis lines=left,
  legend style={font=\footnotesize, draw=layergrey!35, fill=white,
    fill opacity=0.9, text opacity=1, at={(0.015,0.985)},
    anchor=north west, cells={anchor=west}},
  legend cell align=left,
]
\addplot[draw=layerteal!60!black, fill=layerteal!18, fill opacity=0.85,
  line width=0.4pt] coordinates{
  (1,0.0031) (1.25,0.0123) (1.5,0.0309) (1.75,0.0617) (2,0.108) (2.25,0.1728)
  (2.5,0.2469) (2.75,0.321) (3,0.3858) (3.25,0.4321) (3.5,0.4506) (3.75,0.4321)
  (4,0.3858) (4.25,0.321) (4.5,0.2469) (4.75,0.1728) (5,0.108) (5.25,0.0617)
  (5.5,0.0309) (5.75,0.0123) (6,0.0031)};
\addlegendentry{$n=4$}
\addplot[draw=layerindigo!70!black, fill=layerindigo!22, fill opacity=0.8,
  line width=0.4pt] coordinates{
  (1.6875,0) (1.75,0.0001) (1.8125,0.0002) (1.875,0.0004) (1.9375,0.0008) (2,0.0014)
  (2.0625,0.0025) (2.125,0.0043) (2.1875,0.0072) (2.25,0.0117) (2.3125,0.0184) (2.375,0.028)
  (2.4375,0.0416) (2.5,0.0601) (2.5625,0.0847) (2.625,0.1163) (2.6875,0.1558) (2.75,0.2039)
  (2.8125,0.2607) (2.875,0.3258) (2.9375,0.3981) (3,0.4759) (3.0625,0.5567) (3.125,0.6374)
  (3.1875,0.7145) (3.25,0.7843) (3.3125,0.8431) (3.375,0.8878) (3.4375,0.9156) (3.5,0.9251)
  (3.5625,0.9156) (3.625,0.8878) (3.6875,0.8431) (3.75,0.7843) (3.8125,0.7145) (3.875,0.6374)
  (3.9375,0.5567) (4,0.4759) (4.0625,0.3981) (4.125,0.3258) (4.1875,0.2607) (4.25,0.2039)
  (4.3125,0.1558) (4.375,0.1163) (4.4375,0.0847) (4.5,0.0601) (4.5625,0.0416) (4.625,0.028)
  (4.6875,0.0184) (4.75,0.0117) (4.8125,0.0072) (4.875,0.0043) (4.9375,0.0025) (5,0.0014)
  (5.0625,0.0008) (5.125,0.0004) (5.1875,0.0002) (5.25,0.0001) (5.3125,0)};
\addlegendentry{$n=16$}
\addplot[draw=accent!80!black, fill=accent!30, fill opacity=0.8,
  line width=0.5pt] coordinates{
  (2.5312,0) (2.5469,0.0001) (2.5625,0.0001) (2.5781,0.0001) (2.5938,0.0002) (2.6094,0.0003)
  (2.625,0.0004) (2.6406,0.0005) (2.6562,0.0007) (2.6719,0.0009) (2.6875,0.0012) (2.7031,0.0016)
  (2.7188,0.0021) (2.7344,0.0028) (2.75,0.0036) (2.7656,0.0047) (2.7812,0.0061) (2.7969,0.0079)
  (2.8125,0.01) (2.8281,0.0127) (2.8438,0.0161) (2.8594,0.0202) (2.875,0.0252) (2.8906,0.0312)
  (2.9062,0.0385) (2.9219,0.0472) (2.9375,0.0576) (2.9531,0.0698) (2.9688,0.0842) (2.9844,0.1009)
  (3,0.1203) (3.0156,0.1426) (3.0312,0.1681) (3.0469,0.1971) (3.0625,0.2297) (3.0781,0.2663)
  (3.0938,0.307) (3.1094,0.3521) (3.125,0.4015) (3.1406,0.4553) (3.1562,0.5136) (3.1719,0.5762)
  (3.1875,0.6429) (3.2031,0.7134) (3.2188,0.7875) (3.2344,0.8645) (3.25,0.944) (3.2656,1.0252)
  (3.2812,1.1075) (3.2969,1.19) (3.3125,1.2718) (3.3281,1.352) (3.3438,1.4295) (3.3594,1.5035)
  (3.375,1.573) (3.3906,1.6369) (3.4062,1.6943) (3.4219,1.7445) (3.4375,1.7867) (3.4531,1.8202)
  (3.4688,1.8445) (3.4844,1.8592) (3.5,1.8641) (3.5156,1.8592) (3.5312,1.8445) (3.5469,1.8202)
  (3.5625,1.7867) (3.5781,1.7445) (3.5938,1.6943) (3.6094,1.6369) (3.625,1.573) (3.6406,1.5035)
  (3.6562,1.4295) (3.6719,1.352) (3.6875,1.2718) (3.7031,1.19) (3.7188,1.1075) (3.7344,1.0252)
  (3.75,0.944) (3.7656,0.8645) (3.7812,0.7875) (3.7969,0.7134) (3.8125,0.6429) (3.8281,0.5762)
  (3.8438,0.5136) (3.8594,0.4553) (3.875,0.4015) (3.8906,0.3521) (3.9062,0.307) (3.9219,0.2663)
  (3.9375,0.2297) (3.9531,0.1971) (3.9688,0.1681) (3.9844,0.1426) (4,0.1203) (4.0156,0.1009)
  (4.0312,0.0842) (4.0469,0.0698) (4.0625,0.0576) (4.0781,0.0472) (4.0938,0.0385) (4.1094,0.0312)
  (4.125,0.0252) (4.1406,0.0202) (4.1562,0.0161) (4.1719,0.0127) (4.1875,0.01) (4.2031,0.0079)
  (4.2188,0.0061) (4.2344,0.0047) (4.25,0.0036) (4.2656,0.0028) (4.2812,0.0021) (4.2969,0.0016)
  (4.3125,0.0012) (4.3281,0.0009) (4.3438,0.0007) (4.3594,0.0005) (4.375,0.0004) (4.3906,0.0003)
  (4.4062,0.0002) (4.4219,0.0001) (4.4375,0.0001) (4.4531,0.0001) (4.4688,0)};
\addlegendentry{$n=64$}
\draw[layergrey, densely dotted, line width=0.6pt] (axis cs:3.5,0) -- (axis cs:3.5,1.95);
\node[font=\footnotesize, layergrey, anchor=south east] at (axis cs:3.45,1.95) {$3{,}5$};
\end{axis}
\end{tikzpicture}
\caption{Tiga sebaran rata-rata $n$ lemparan dadu, digambar pada satu skala yang
sama: $n=4$, $16$, dan $64$. Lebarnya $0{,}853913$, $0{,}426956$, dan
$0{,}213478$, jadi setiap kali $n$ dilipat empatkan lebarnya tinggal separuh.
Puncaknya meninggi dari $0{,}4506$ ke $0{,}9251$ dan $1{,}8641$, sedangkan
pusatnya tetap di $3{,}5$.}
\label{fig:tiga-lebar}
\end{figure}

Angka-angka membuat perbedaan itu terasa. Untuk lemparan koin dengan $p=1/2$ dan
toleransi $\varepsilon=0{,}1$, Chebyshev memberi
$\sigma^2/(n\varepsilon^2)=0{,}25/(0{,}01n)=25/n$. Agar peluang meleset paling
besar $0{,}3$, cukup $n\ge25/0{,}3=83{,}33$, yaitu $n=84$. Pada $n=84$ batas itu
bernilai $0{,}297619$, jadi Chebyshev menjamin peluang berada dalam
$0{,}4$ sampai $0{,}6$ paling sedikit $0{,}702381$. Peluang sebenarnya, dihitung
dari sebaran binomial, adalah $0{,}937028$, dan aproksimasi normal memberi
$0{,}933197$. Jaminan itu benar tetapi terlalu rendah sekitar $23$ poin
persentase.

Kelemahan yang sama menjadi parah bila $\sigma$ besar. Untuk dadu adil,
$\operatorname{Var}(X)=35/12=2{,}916667$ dan $\sigma=1{,}707825$. Pada $n=100$
dengan $\varepsilon=0{,}1$, Chebyshev memberi batas $2{,}916667$, angka yang lebih
besar dari satu dan karena itu tidak berarti apa pun. Untuk mendapatkan galat
$0{,}01$ dengan keyakinan $95$ persen, rumus normal menuntut paling sedikit
$112043$ pengamatan, sedangkan Chebyshev menuntut jauh lebih banyak lagi.

Di sinilah letak persoalan yang menyiapkan bagian berikutnya. Chebyshev menjawab
"apakah rata-rata boleh dipercaya" dengan benar tetapi tumpul. Ia tidak mampu
menjawab "seberapa dekat, dengan peluang berapa" secara layak dipakai. Jawaban
tajam menuntut pengetahuan tentang bentuk sebaran galat, bukan hanya ukurannya.

\begin{bacakeras}
Dua suara, bergantian. Yang pertama menyebut banyaknya lemparan; yang kedua
menyebut jarak meleset yang masih wajar.

--- Sepuluh lemparan. Melesetnya masih lebar.
--- Seratus lemparan. Sudah lebih sempit.
--- Seribu lemparan. Makin sempit lagi.

Sekarang lipat empatkan, jangan tambah sepuluh: sepuluh jadi empat puluh,
seratus jadi empat ratus, seribu jadi empat ribu. Melesetnya tinggal separuh.
Bukan seperempat. Separuh. Setiap kali empat kali, tepat separuh.

Lalu balik arahnya, dan bacakan lebih pelan. Berapa ongkos satu angka tambahan
di belakang koma? Seratus kali lemparan. Seratus kali kerja, satu angka.
Rata-rata tidak sabar: ia menuntut empat kali, bukan dua kali, dan tidak
menawarkan diskon.

Baris yang diulang tiga kali: empat kali data, separuh galat.
\end{bacakeras}

\begin{untukanak}
Bayangkan mesin pengisi kantong beras. Setiap kantong meleset sedikit dari berat
yang diminta, kadang lebih kadang kurang, dan besar melesetnya berbeda-beda;
menimbang satu kantong tidak memberi tahu arah melesetnya.

Sekarang mesin itu mengisi seratus kantong, lalu kita menimbang semuanya
sekaligus. Meleset ke atas dan ke bawah saling memotong, sehingga berat totalnya
jauh lebih dekat ke seratus kali berat yang diminta. Kalau kantongnya sepuluh
ribu, jumlahnya hampir tepat.

Yang mengejutkan, bentuk sebaran meleset itu tidak perlu kita ketahui. Entah
melesetnya berupa lonceng, berupa dua nilai saja, atau berupa apa pun yang aneh,
sebaran jumlahnya tetap menyerupai lonceng yang sama, asalkan kantongnya cukup
banyak dan melesetnya tidak terlalu sering menjadi raksasa. Inilah yang membuat
para penjudi, pembuat jembatan, dan peneliti cuaca bisa memakai rumus yang sama.

Pertanyaan yang membuat penasaran: mengapa bentuk asli melesetnya tidak lagi
penting setelah kantongnya banyak? Dan mengapa satu kantong yang melesetnya
raksasa bisa merusak seluruh kesimpulan? Di bab ini kita akan tahu mengapa.
\end{untukanak}

\section{Hukum yang kuat: apa yang sesungguhnya ditambahkan}
\tingkat{lanjut}

Kalimat sehari-harinya: dengan keyakinan satu, rata-rata pada satu jalankan
panjang akhirnya menetap, bukan sekadar sering dekat.

Hukum lemah berbicara tentang setiap $n$ secara terpisah: pada $n$ yang cukup
besar, peluang meleset kecil. Ia tidak melarang rata-rata meleset jauh pada
$n=10^6$, lalu kembali tenang, lalu meleset lagi pada $n=10^9$, asalkan kejadian
semacam itu jarang pada setiap $n$. Hukum kuat menutup celah itu.

\begin{teorema}[Hukum bilangan besar kuat]
Misalkan $X_1,X_2,\dots$ bebas dan berdistribusi sama dengan
$\mathbb{E}|X_1|<\infty$ dan $\mathbb{E}[X_1]=\mu$. Maka
$\bar X_n\to\mu$ hampir pasti.\index{hukum bilangan besar kuat}
\end{teorema}

Bedanya halus tetapi tajam. Barisan $X_n$ pada bagian sebelumnya konvergen dalam
peluang ke nol, tetapi pada setiap $\omega$ nilainya satu tak hingga sering.
Hukum lemah tidak dapat menolaknya; hukum kuat menolaknya, karena ia menuntut
konvergensi pada setiap $\omega$ di luar himpunan berpeluang nol.

Perhatikan juga syaratnya. Hukum kuat hanya menuntut $\mathbb{E}|X_1|<\infty$,
tanpa variansi. Hukum lemah yang kita buktikan di atas menuntut $\sigma^2<\infty$
hanya karena buktinya memakai Chebyshev; syarat itu bukan syarat alami teoremanya.
Untuk hukum lemah, syarat perlu dan cukup bagi barisan iid adalah
$n\,P(|X_1|>n)\to0$, yang lebih lemah daripada $\mathbb{E}|X_1|<\infty$. Rincian
itu kita tinggalkan, dan pembaca dapat menemukannya pada Feller, \emph{An
Introduction to Probability Theory and Its Applications}, jilid 2, bab tentang
hukum bilangan besar.

Bukti hukum kuat tidak kita kerjakan di sini. Bentuk klasiknya memotong variabel
pada ambang $n$, menerapkan Borel--Cantelli pada ekor potongan, lalu memakai
kriteria Cauchy untuk kekonvergenan hampir pasti; Etemadi memberi bukti yang
menghindari asumsi variansi dan hanya memakai pasangan variabel yang saling
bebas. Lokasi buktinya: Blitzstein dan Hwang, \emph{Introduction to Probability},
bab 10 bagian 10.2; Wasserman, \emph{All of Statistics}, lampiran bab 5 bagian
5.7.1. Yang perlu diingat dari pernyataan itu: hukum kuat tidak memberi laju.
Ia mengatakan bahwa galat akhirnya kecil dan tinggal kecil, bukan seberapa cepat.

\begin{tangan}
Ambil satu koin dan selembar kertas bergaris. Angka bernilai satu, gambar
bernilai nol.

Lempar empat puluh kali, tulis 1 atau 0 berderet. Di baris kedua, tulis
rata-rata berjalan di bawah setiap lemparan: jumlah sampai situ dibagi
banyaknya lemparan sampai situ.

Yang harus kita lihat: baris kedua bergoyang paling liar pada lima lemparan
pertama, lalu makin tenang. Pada lemparan keempat puluh angkanya belum tentu
$0{,}5$; masih wajar $0{,}475$ atau $0{,}525$.

Sekarang ulangi percobaan itu empat kali dan gambar keempat jalur di satu
grafik dengan pensil berbeda. Keempat jalur menyempit ke $0{,}5$, dan lebar
sebarannya tinggal separuh setiap kali lemparannya dilipat empatkan.

Tempel kertasnya di dinding. Kalau ada yang berkata hasilnya pasti seimbang,
tunjuk lemparan keempat puluh, bukan yang kesepuluh.
\end{tangan}

\section{Teorema limit pusat}
\tingkat{lanjut}

Kalimat sehari-harinya: galat rata-rata, setelah dibesarkan dengan faktor
$\sqrt n$, selalu berbentuk lonceng yang sama, apa pun bentuk asli datanya.

Hukum bilangan besar memberi tahu bahwa $\bar X_n$ menempel ke $\mu$. Ia tidak
memberi tahu seberapa lebar tempelan itu. Teorema limit pusat menjawabnya, dan
jawabannya mengejutkan karena tidak bergantung pada distribusi asal.

\kamus{teorema limit pusat}{central limit theorem}{bentuk lonceng yang muncul pada
galat rata-rata, bukan pada datanya, dan hanya setelah galat itu dibesarkan
sebesar akar banyaknya data}

\begin{teorema}[Limit pusat, bentuk Lindeberg--Levy]\index{teorema limit pusat}
Misalkan $X_1,X_2,\dots$ bebas dan berdistribusi sama dengan
$\mathbb{E}[X_i]=\mu$ dan $0<\operatorname{Var}(X_i)=\sigma^2<\infty$. Maka
\[
Z_n=\frac{\sqrt n\,(\bar X_n-\mu)}{\sigma}
\;\rightsquigarrow\; Z\sim N(0,1),
\]
artinya $P(Z_n\le z)\to\Phi(z)$ untuk setiap $z\in\mathbb{R}$.
\end{teorema}

Perhatikan letak $\sqrt n$. Simpangan baku $\bar X_n$ adalah
$\sigma/\sqrt n$, bukan $\sigma/n$. Karena itu untuk memangkas galat menjadi
separuh kita perlu empat kali data, dan untuk menambah satu digit ketelitian kita
perlu seratus kali data. Ini kabar buruk yang tidak bisa diakali: ketelitian
datang dengan harga kuadratik.

Gambar~\ref{fig:lonceng-dari-kotak} memperlihatkan dari mana bentuk lonceng itu
datang. Panel kiri adalah sebaran satu dadu: datar, tanpa lonceng sama sekali.
Panel kanan adalah sebaran rata-rata dua belas dadu yang sudah distandarkan.
Bahannya bukan lonceng, hasilnya lonceng.

\begin{figure}[htbp]
\centering
\begin{tikzpicture}
\begin{axis}[
  name=satu, width=4.3cm, height=4.5cm,
  ybar, bar width=11pt,
  xmin=0.5, xmax=6.5, ymin=0, ymax=0.21,
  xtick={1,2,3,4,5,6}, ytick={0,0.0833,0.1667},
  yticklabels={$0$,$\frac{1}{12}$,$\frac16$},
  xlabel={mata dadu}, ylabel={peluang},
  title={\footnotesize satu dadu},
  label style={font=\footnotesize}, tick label style={font=\scriptsize},
  axis lines=left,
]
\addplot[fill=layerindigo!50, draw=layerindigo!80!black, line width=0.3pt]
  coordinates{
  (1,0.1667) (2,0.1667) (3,0.1667) (4,0.1667) (5,0.1667) (6,0.1667)};
\end{axis}
\begin{axis}[
  name=duabelas, at={(satu.east)}, anchor=west, xshift=7pt,
  width=9.3cm, height=4.5cm,
  ybar, bar width=3.4pt,
  xmin=-5.3, xmax=5.3, ymin=0, ymax=0.45,
  xtick={-4,-2,0,2,4}, ytick={0,0.2,0.4},
  xlabel={skor $z$ dari rata-rata}, ylabel={kerapatan},
  title={\footnotesize dua belas dadu, sudah distandarkan},
  label style={font=\footnotesize}, tick label style={font=\scriptsize},
  axis lines=left,
]
\addplot[fill=accent!45, draw=accent!70!black, line width=0.2pt]
  coordinates{
  (-4.3948,0) (-4.2258,0.00001) (-4.0567,0.00003) (-3.8877,0.00009) (-3.7187,0.0002) (-3.5496,0.00044)
  (-3.3806,0.00091) (-3.2116,0.00177) (-3.0426,0.00327) (-2.8735,0.00575) (-2.7045,0.00966) (-2.5355,0.01559)
  (-2.3664,0.02418) (-2.1974,0.03614) (-2.0284,0.05215) (-1.8593,0.07276) (-1.6903,0.09826) (-1.5213,0.1286)
  (-1.3522,0.16323) (-1.1832,0.20113) (-1.0142,0.24073) (-0.8452,0.28001) (-0.6761,0.31669) (-0.5071,0.34838)
  (-0.3381,0.37287) (-0.169,0.38835) (0,0.39365) (0.169,0.38835) (0.3381,0.37287) (0.5071,0.34838)
  (0.6761,0.31669) (0.8452,0.28001) (1.0142,0.24073) (1.1832,0.20113) (1.3522,0.16323) (1.5213,0.1286)
  (1.6903,0.09826) (1.8593,0.07276) (2.0284,0.05215) (2.1974,0.03614) (2.3664,0.02418) (2.5355,0.01559)
  (2.7045,0.00966) (2.8735,0.00575) (3.0426,0.00327) (3.2116,0.00177) (3.3806,0.00091) (3.5496,0.00044)
  (3.7187,0.0002) (3.8877,0.00009) (4.0567,0.00003) (4.2258,0.00001) (4.3948,0)};
\addplot[domain=-4.6:4.6, samples=120, layerindigo, line width=0.9pt]
  {0.39894228*exp(-x^2/2)};
\end{axis}
\end{tikzpicture}
\caption{Dari datar ke lonceng. Kiri: sebaran satu dadu, rata di enam nilai.
Kanan: sebaran rata-rata dua belas dadu setelah distandarkan, ditumpuk dengan
lengkung normal baku. Jarak supremum keduanya $0{,}03384$, kurang dari separuh
jarak pada $n=2$ yang $0{,}08333$. Bahan asalnya tidak lonceng, hasilnya
lonceng.}
\label{fig:lonceng-dari-kotak}
\end{figure}

Bukti berikut memakai fungsi pembangkit momen. Ia lengkap untuk kasus yang momen
ketiganya hingga, dan kita nyatakan batasnya secara jujur di akhir.

\begin{lemma}[Kekonvergenan lewat fungsi pembangkit momen]\index{fungsi pembangkit momen}
Misalkan $\psi_n(t)=\mathbb{E}[e^{tZ_n}]$ dan $\psi(t)=\mathbb{E}[e^{tZ}]$ hingga
pada suatu selang terbuka di sekitar $t=0$. Jika $\psi_n(t)\to\psi(t)$ untuk
setiap $t$ di selang itu, maka $Z_n\rightsquigarrow Z$.
\end{lemma}

Lemma itu adalah teorema kontinuitas untuk fungsi pembangkit momen. Buktinya
memakai kenyataan bahwa fungsi sebaran ditentukan oleh fungsi pembangkit momennya
pada selang di sekitar nol, dan kita tidak mengulanginya di sini; lihat Wasserman,
\emph{All of Statistics}, bab 5 bagian 5.7.2, Lemma 5.20.

\begin{proof}[Bukti teorema limit pusat]
Tetapkan $Y_i=(X_i-\mu)/\sigma$, sehingga $\mathbb{E}[Y_i]=0$,
$\operatorname{Var}(Y_i)=1$, dan $Z_n=n^{-1/2}\sum_{i=1}^n Y_i$. Misalkan
$\psi(t)=\mathbb{E}[e^{tY_1}]$ fungsi pembangkit momen $Y_1$, yang hingga pada
selang $|t|\le\delta$. Karena variabelnya bebas, fungsi pembangkit momen $Z_n$
adalah
\[
\xi_n(t)=\Bigl[\psi\bigl(t/\sqrt n\bigr)\Bigr]^n .
\]
Turunan di nol memberi $\psi(0)=1$, $\psi'(0)=\mathbb{E}[Y_1]=0$, dan
$\psi''(0)=\mathbb{E}[Y_1^2]=1$. Ekspansi Taylor orde dua dengan sisa orde tiga,
yang sah bila $\mathbb{E}|Y_1|^3<\infty$, berbunyi
\[
\psi(s)=1+\frac{s^2}{2}+R(s),\qquad |R(s)|\le C|s|^3 \quad (|s|\le\delta).
\]
Untuk $n$ cukup besar berlaku $|t|/\sqrt n\le\delta$, sehingga
\[
\xi_n(t)=\Bigl(1+\frac{t^2}{2n}+R\bigl(t/\sqrt n\bigr)\Bigr)^n
=\Bigl(1+\frac{t^2}{2n}+O\bigl(n^{-3/2}\bigr)\Bigr)^n .
\]
Ambil logaritma. Dengan $\log(1+a)=a-a^2/2+\dots$ dan
$a_n=t^2/(2n)+O(n^{-3/2})$,
\[
\log\xi_n(t)=n\Bigl(a_n-\frac{a_n^2}{2}+\dots\Bigr)
=\frac{t^2}{2}+O\bigl(n^{-1/2}\bigr)-\frac{t^4}{8n}+\dots
\longrightarrow \frac{t^2}{2}.
\]
Jadi $\xi_n(t)\to e^{t^2/2}$ untuk setiap $t$ tetap. Fungsi terakhir adalah fungsi
pembangkit momen $N(0,1)$, sehingga Lemma di atas memberi
$Z_n\rightsquigarrow N(0,1)$.
\end{proof}

Ada satu syarat tambahan yang menyelinap masuk: sisa Taylor orde tiga menuntut
$\mathbb{E}|Y_1|^3<\infty$. Untuk distribusi berekor sangat berat syarat itu bisa
gagal walaupun variansinya hingga, dan bukti di atas harus diganti. Jalan yang
umum adalah memakai fungsi karakteristik $\varphi(t)=\mathbb{E}[e^{itX}]$, yang
selalu ada karena $|e^{itX}|=1$, lalu memakai teorema kontinuitas Levy. Rangka
buktinya sama: logaritma fungsi karakteristik $Z_n$ menuju $-t^2/2$. Versi tanpa
asumsi momen ketiga ada di Wasserman, bab 5 bagian 5.7.2, dan di Blitzstein dan
Hwang, bab 10 bagian 10.3.

Teorema ini tidak mengatakan bahwa galatnya kecil. Ia mengatakan bahwa galat itu
berbentuk lonceng setelah dibagi $\sigma/\sqrt n$. Ukurannya tetap diberikan oleh
$\sigma/\sqrt n$, dan $\sigma$ adalah besaran yang biasanya tidak kita ketahui.
Dalam praktik $\sigma$ diganti oleh simpangan baku sampel $s$, dan penggantian itu
tidak merusak limit; pernyataannya menjadi $\sqrt n(\bar X_n-\mu)/s_n
\rightsquigarrow N(0,1)$. Besaran $\sigma/\sqrt n$ disebut galat baku
\index{galat baku} (\emph{standard error}), dan inilah dasar interval kepercayaan
di Bab 20.

Seberapa cepat aproksimasi itu membaik juga bisa dibatasi.

\begin{teorema}[Ketaksamaan Berry--Esseen]
Jika $\rho=\mathbb{E}|X_1-\mu|^3<\infty$, maka
\[
\sup_z\bigl|P(Z_n\le z)-\Phi(z)\bigr|\le \frac{C\rho}{\sigma^3\sqrt n},
\]
dengan $C$ konstanta universal; nilai $C=0{,}4748$ cukup untuk semua distribusi.
\end{teorema}

Bentuk ketaksamaan ini penting: jaraknya mengecil seperti $1/\sqrt n$, dengan
konstanta yang bergantung pada kemencengan distribusi asal lewat $\rho/\sigma^3$.
Distribusi yang sangat menceng akan butuh $n$ jauh lebih besar daripada
distribusi yang simetris. Angka pada bagian lab menunjukkan hal ini secara
kuantitatif untuk distribusi eksponensial, yang momen ketiganya $\rho=2{,}414553$
dan karenanya batasnya $1{,}1464/\sqrt n$.

\begin{resep}
Tiga langkah menormalkan rata-rata sampel. Urutannya tidak boleh ditukar.
\begin{enumerate}
\item \textbf{Kurangi pusat.} Hitung $\bar x-\mu$. Lonceng bergeser sehingga
      puncaknya jatuh di nol.
\item \textbf{Bagi lebar.} Bagi hasilnya dengan galat baku $\sigma/\sqrt n$.
      Bentuk lonceng tidak berubah; hanya skalanya yang menjadi satu satuan.
\item \textbf{Baca tabel.} Sisa angka $z$ dibaca pada $\Phi(z)$; untuk ekor
      kanan pakai $1-\Phi(z)$.
\end{enumerate}
Peringatan urutan: membagi sebelum mengurangi pusat menggeser hasil sebesar
$\mu\sqrt n/\sigma$. Pada dadu dengan $\mu=3{,}5$, $\sigma=1{,}707825$, dan
$n=12$, pergeseran itu $7{,}10$; angkanya sudah di luar tabel normal mana pun.
\end{resep}

\section{Aproksimasi normal dan koreksi kontinuitas}
\tingkat{menengah}

Kalimat sehari-harinya: untuk menghitung peluang jumlah, kita boleh memakai
tabel normal, tetapi bila jumlahnya hanya bisa bernilai bulat, ambangnya digeser
setengah.

Aproksimasi normal untuk jumlah bekerja paling sederhana pada distribusi diskret,
dan di situ muncul satu koreksi kecil yang dampaknya besar. Namanya koreksi
kontinuitas\index{koreksi kontinuitas}. Ambil $S\sim
\text{Binomial}(100;\,0{,}5)$, sehingga $\mu=50$ dan $\sigma=\sqrt{100\cdot0{,}5
\cdot0{,}5}=5$. Kejadian $\{S\le40\}$ adalah kejadian yang sama dengan
$\{S<40{,}5\}$, karena $S$ hanya bernilai bulat. Karena itu ambang yang benar
untuk distandarkan adalah $40{,}5$, bukan $40$.

\begin{center}
\begin{tabular}{lrrr}
\hline
Kejadian & Eksak & Tanpa koreksi & Dengan koreksi \\
\hline
$P(S\le40)$ & $0{,}028444$ & $0{,}022750$ & $0{,}028717$ \\
$P(S\le55)$ & $0{,}864373$ & $0{,}841345$ & $0{,}864334$ \\
$P(S=60)$   & $0{,}010844$ & $0{,}010798$ & $0{,}010852$ \\
\hline
\end{tabular}
\end{center}

Bacalah tabel itu baris demi baris. Pada baris pertama, aproksimasi tanpa koreksi
meleset $0{,}005694$, sedangkan dengan koreksi meleset $0{,}000273$, hampir dua
puluh satu kali lebih dekat. Pada baris ketiga keduanya sudah dekat, karena di
sekitar pusat sebaran lebar satu satuan sudah kecil dibandingkan $\sigma=5$.
Pelajaran umumnya: koreksi kontinuitas paling menolong di ekor dan ketika
$\sigma$ kecil. Besar pergeseran $z$ yang ditimbulkannya adalah $1/(2\sigma)$,
yaitu $0{,}1$ pada contoh ini.

\kertas{Hitung ulang baris pertama tabel ini dengan tangan. $z=(40-50)/5=-2$,
$\Phi(-2)=0{,}022750$; lalu $z=(40{,}5-50)/5=-1{,}9$, $\Phi(-1{,}9)=0{,}028717$.
Nilai eksaknya $\sum_{k=0}^{40}\binom{100}{k}2^{-100}=0{,}028444$. Periksa juga
bahwa galat tanpa koreksi $0{,}005694$ sekitar $21$ kali galat dengan koreksi
$0{,}000273$.}

Koreksi ini bukan rumus baru, hanya pembacaan ulang kejadian. Untuk variabel
diskret dengan jarak antar-nilai $h$, ambang digeser $h/2$ ke arah dalam. Untuk
variabel kontinu tidak ada yang perlu digeser, karena tidak ada lompatan.

Satu kebiasaan buruk perlu disebut di sini. Banyak buku menulis bahwa
aproksimasi normal "boleh dipakai bila $n\ge30$". Itu bukan teorema melainkan
kebiasaan. Angka pada bagian lab memperlihatkan mengapa: untuk rata-rata
eksponensial, jarak supremum ke normal pada $n=30$ adalah $0{,}02429$, dan
kuantil $95$ persen dari $Z_{30}$ adalah $1{,}7419$ sedangkan normal memberi
$1{,}6449$. Selisih $0{,}097$ itu berarti interval kepercayaan yang terlalu
sempit sekitar enam persen. Yang menentukan bukan $n$ saja, melainkan kemencengan
dan kurtosis distribusi asal.

\section{Bila teorema limit pusat menyesatkan: ekor berat}
\tingkat{lanjut}

Kalimat sehari-harinya: bila satu pengamatan sesekali bisa jauh lebih besar
daripada semua yang lain, rata-rata tidak lagi menetap, dan lonceng tidak muncul.

\kamus{ekor berat}{heavy tail}{peluang nilai raksasa yang tidak cukup cepat
mengecil, sehingga satu pengamatan dapat mengalahkan seluruh sisanya}

Semua teorema di atas memakai satu syarat yang mudah dilewati: variansi hingga.
Distribusi Cauchy\index{distribusi Cauchy}, dengan kepadatan
\[
f(x)=\frac{1}{\pi(1+x^2)},\qquad x\in\mathbb{R},
\]
tidak memilikinya. Ekornya mengecil seperti $1/x$, terlalu lambat. Peluang
melewati sepuluh adalah $1-(2/\pi)\arctan 10=0{,}063451$, dan peluang melewati
seratus masih $0{,}006366$. Bandingkan dengan normal baku, yang peluang melewati
sepuluh praktis nol. Bahkan $\mathbb{E}|X|$ tidak hingga: harapan terpangkas pada
ambang $M$ bernilai $(2/\pi)\ln(1+M^2)$, yaitu $2{,}9381$ untuk $M=10$,
$8{,}7952$ untuk $M=10^3$, dan $17{,}5905$ untuk $M=10^6$. Angka itu tumbuh
seperti logaritma, jadi tidak ada batas.

Akibatnya tegas, dan bukan sekadar "CLT kurang akurat".

\begin{proposisi}
Jika $X_1,\dots,X_n$ iid Cauchy, maka $\bar X_n$ berdistribusi Cauchy untuk
setiap $n$.
\end{proposisi}

\begin{proof}[Bukti]
Fungsi karakteristik Cauchy adalah $\varphi(t)=e^{-|t|}$. Karena variabelnya
bebas, fungsi karakteristik $\bar X_n$ adalah
$\bigl[e^{-|t|/n}\bigr]^n=e^{-|t|}$, yaitu fungsi karakteristik Cauchy lagi.
Karena fungsi karakteristik menentukan sebaran, $\bar X_n$ berdistribusi Cauchy.
\end{proof}

Jadi rata-rata $10^5$ pengamatan Cauchy sama tidak informatifnya dengan satu
pengamatan. Hukum bilangan besar gagal, teorema limit pusat gagal, dan "galat
baku" tidak berarti apa pun. Angka pada bagian lab menunjukkan tiga gejala itu
sekaligus. Median $|\bar X_n|$ tinggal di sekitar satu untuk semua $n$
($1{,}004$, $1{,}008$, $1{,}040$, $0{,}973$), padahal untuk distribusi
bervariansi hingga ia harus mengecil seperti $1/\sqrt n$. Peluang
$|\bar X_n|>10$ tinggal di sekitar $0{,}0635$, nilai teoretisnya, untuk semua $n$.
Simpangan baku empirisnya melompat-lompat: $33{,}04$, $863{,}61$, $23{,}87$,
$379{,}53$ untuk $n=10$ sampai $10^4$. Lompatan itu bukan kesalahan program.
Simpangan baku sampel dari data Cauchy sendiri adalah variabel acak yang tidak
punya nilai khas.

Gambar~\ref{fig:cauchy-berjalan} memperlihatkan bedanya pada satu jalankan.
Rata-rata berjalan data seragam menempel ke $0{,}5$ sesudah beberapa puluh
tarikan. Rata-rata berjalan data Cauchy melompat dari $-0{,}0866$ ke $4{,}4572$
pada tarikan ke-20, karena satu nilai $90{,}79$ kebetulan masuk, lalu mengembara
lagi. Satu nilai raksasa cukup untuk membatalkan seluruh riwayat sebelumnya.

\begin{figure}[htbp]
\centering
\begin{tikzpicture}
\begin{axis}[
  width=0.95\linewidth, height=6.0cm,
  xmode=log, xmin=1, xmax=4000, ymin=-1.1, ymax=4.9,
  xlabel={banyak tarikan $n$ (skala logaritma)}, ylabel={rata-rata berjalan},
  label style={font=\small}, tick label style={font=\footnotesize},
  axis lines=left,
  legend style={font=\footnotesize, draw=layergrey!35, fill=white,
    fill opacity=0.9, text opacity=1, at={(0.015,0.985)},
    anchor=north west, cells={anchor=west}},
  legend cell align=left,
]
\addplot[accent, line width=0.7pt] coordinates{
  (1,-0.253) (2,1.045) (3,0.582) (4,0.564) (5,0.553) (6,1.064)
  (7,-0.579) (8,-0.517) (9,-0.409) (10,-0.325) (11,-0.409) (12,-0.361)
  (13,-0.344) (14,-0.262) (15,0.008) (16,-0.038) (17,0.017) (18,-0.08)
  (19,-0.087) (20,4.457) (21,4.29) (22,4.153) (23,4.046) (24,3.762)
  (25,3.656) (26,3.551) (27,3.425) (28,3.257) (29,3.186) (30,3.083)
  (31,2.934) (32,2.786) (33,2.837) (34,2.743) (35,2.633) (36,2.528)
  (37,2.464) (38,2.381) (39,2.279) (40,2.218) (41,2.157) (44,1.732)
  (47,1.387) (50,1.367) (53,1.191) (57,1.113) (61,1.077) (65,0.942)
  (70,0.796) (75,0.773) (80,0.725) (85,0.492) (91,0.79) (97,0.655)
  (104,0.54) (111,0.478) (119,0.494) (127,0.401) (135,0.981) (145,0.868)
  (155,2.5) (165,2.351) (177,2.207) (189,2.168) (202,1.935) (216,1.854)
  (230,1.776) (246,1.665) (263,1.573) (281,0.848) (300,0.794) (321,0.791)
  (343,0.511) (367,0.922) (392,0.986) (419,0.984) (447,1.03) (478,0.951)
  (511,0.827) (546,0.685) (583,0.37) (623,0.548) (666,0.953) (712,1.021)
  (761,0.945) (813,2.066) (869,2.135) (929,1.855) (992,1.814) (1060,1.649)
  (1133,1.31) (1211,0.557) (1294,0.422) (1383,0.434) (1478,0.466) (1579,0.324)
  (1688,0.406) (1803,0.386) (1927,0.444) (2059,0.403) (2201,0.585) (2352,0.419)
  (2513,0.528) (2686,0.397) (2870,0.356) (3067,0.326) (3278,0.276) (3503,-0.096)
  (3743,-0.159) (4000,-0.203)};
\addlegendentry{Cauchy}
\addplot[layerteal, dashed, line width=1.0pt] coordinates{
  (1,0.221) (2,0.272) (3,0.323) (4,0.472) (5,0.462) (6,0.519)
  (7,0.524) (8,0.573) (9,0.55) (10,0.574) (11,0.605) (12,0.601)
  (13,0.57) (14,0.581) (15,0.543) (16,0.532) (17,0.529) (18,0.528)
  (19,0.532) (20,0.544) (21,0.545) (22,0.547) (23,0.539) (24,0.531)
  (25,0.54) (26,0.555) (27,0.561) (28,0.541) (29,0.545) (30,0.547)
  (31,0.562) (32,0.549) (33,0.533) (34,0.542) (35,0.531) (36,0.535)
  (37,0.537) (38,0.526) (39,0.519) (40,0.516) (41,0.517) (44,0.518)
  (47,0.519) (50,0.503) (53,0.518) (57,0.515) (61,0.525) (65,0.53)
  (70,0.529) (75,0.53) (80,0.535) (85,0.538) (91,0.539) (97,0.538)
  (104,0.541) (111,0.535) (119,0.532) (127,0.533) (135,0.528) (145,0.526)
  (155,0.518) (165,0.523) (177,0.513) (189,0.516) (202,0.513) (216,0.507)
  (230,0.5) (246,0.498) (263,0.495) (281,0.496) (300,0.496) (321,0.495)
  (343,0.499) (367,0.499) (392,0.501) (419,0.503) (447,0.497) (478,0.505)
  (511,0.508) (546,0.509) (583,0.504) (623,0.502) (666,0.497) (712,0.494)
  (761,0.488) (813,0.484) (869,0.485) (929,0.484) (992,0.485) (1060,0.484)
  (1133,0.483) (1211,0.485) (1294,0.487) (1383,0.489) (1478,0.488) (1579,0.491)
  (1688,0.489) (1803,0.489) (1927,0.49) (2059,0.491) (2201,0.49) (2352,0.492)
  (2513,0.491) (2686,0.493) (2870,0.492) (3067,0.49) (3278,0.492) (3503,0.491)
  (3743,0.493) (4000,0.493)};
\addlegendentry{seragam $(0,1)$}
\draw[layergrey, densely dotted, line width=0.6pt]
  (axis cs:1,0.5) -- (axis cs:4000,0.5);
\node[font=\scriptsize, layergrey, anchor=south west] at (axis cs:1.3,0.53) {$0{,}5$};
\end{axis}
\end{tikzpicture}
\caption{Rata-rata berjalan dari 4000 tarikan. Garis putus-putus: seragam pada
$(0,1)$, menetap dekat $0{,}5$ dan berakhir di $0{,}4930$. Garis penuh: Cauchy,
melompat dari $-0{,}0866$ ke $4{,}4572$ pada tarikan ke-20 ketika satu nilai
$90{,}79$ masuk, lalu mengembara lagi dan berakhir di $-0{,}203$. Lebar sumbu
mendatar logaritma, supaya lompatan awal tidak tertutup.}
\label{fig:cauchy-berjalan}
\end{figure}

Yang menggantikan teorema limit pusat adalah keluarga sebaran stabil
\index{sebaran stabil}. Untuk
variabel iid dengan ekor yang mengecil seperti $x^{-\alpha}$ dengan
$0<\alpha<2$, ada konstanta $a_n$ sehingga
\[
\frac{S_n-a_n}{n^{1/\alpha}}\;\rightsquigarrow\;Y_\alpha,
\]
dengan $Y_\alpha$ sebaran stabil berindeks $\alpha$. Untuk $\alpha=1$ kita
mendapat Cauchy dan $n^{1/1}=n$, persis seperti proposisi di atas. Untuk
$\alpha\in(1,2)$ rata-rata masih konvergen ke $\mu$, tetapi galatnya mengecil
seperti $n^{1-1/\alpha}$, lebih lambat daripada $1/\sqrt n$, dan bentuknya bukan
lonceng. Pernyataan umum ini adalah hasil Gnedenko dan Kolmogorov; rinciannya ada
pada Feller jilid 2, bab tentang sebaran stabil.

Bila $\alpha>2$ tetapi momen ketiga tidak ada, teorema limit pusat tetap benar
karena variansinya hingga. Yang hilang adalah jaminan lajunya: batas
Berry--Esseen menjadi tak berhingga, sehingga tidak ada lagi pegangan bahwa galat
mengecil seperti $1/\sqrt n$ dengan konstanta yang terkendali. Dalam praktik ini
yang paling sering menggigit: sebarannya tidak ekstrem seperti Cauchy, tetapi
cukup berekor untuk membuat interval kepercayaan terlalu sempit pada ukuran
sampel yang realistis.

\begin{tantangan}
Distribusi tanpa rata-rata. Ambil $X$ dengan
$P(X=2^k)=P(X=-2^k)=2^{-k-1}$ untuk $k=1,2,3,\dots$

Tunjukkan dua hal. Pertama, $\mathbb{E}|X|=\infty$, sehingga tidak ada $\mu$
yang boleh dituju. Kedua, rata-rata sampelnya tidak konvergen dalam peluang ke
konstanta mana pun: hukum bilangan besar gagal, bukan sekadar melambat.

Petunjuk 1: yang menentukan bukan rata-rata semua suku, melainkan suku
terbesar. Tanyakan seberapa sering satu suku melewati ambang $n$.

Petunjuk 2: hitung $n\,P(|X|>n)$ untuk $n$ di antara dua pangkat dua, lalu
pakai syarat perlu $n\,P(|X_1|>n)\to0$ yang disebut pada bagian hukum kuat.

Kalau mentok, kembali ke bagian tentang ekor berat dan ke paragraf syarat perlu
dan cukup di bagian hukum kuat.
\end{tantangan}

\begin{untukanak}
Di sebuah desa ada seratus orang dengan penghasilan biasa-biasa saja, dan satu
orang yang sangat kaya. Kalau kita menghitung penghasilan rata-rata desa itu,
angkanya ditentukan hampir sepenuhnya oleh orang kaya tadi. Kalau orang kaya itu
kebetulan sedang pergi, rata-ratanya jatuh drastis. Kalau ia pulang, rata-ratanya
melompat naik.

Sekarang tambahkan penduduk. Seratus orang lagi, seribu orang lagi; selama
pendatangnya orang biasa, rata-ratanya berubah sedikit. Tetapi kalau di antara
pendatang itu ada satu orang raksasa lagi, rata-ratanya melompat lagi. Menambah
banyak orang tidak menolong, karena yang menentukan bukan banyaknya orang,
melainkan siapa yang terbesar.

Inilah yang terjadi pada data yang ekornya berat: satu pengamatan bisa
mengalahkan semua yang lain. Rata-rata menjadi tidak stabil, dan lonceng yang
biasa dipakai untuk menghitung "seberapa yakin" tidak muncul sama sekali.

Pertanyaan yang membuat penasaran: bagaimana kita bisa tahu, sebelum menghitung,
apakah data kita termasuk jenis yang tenang atau jenis yang dikuasai raksasa? Di
bab ini kita akan tahu mengapa pertanyaan itu tidak bisa dijawab hanya dengan
melihat rata-ratanya.
\end{untukanak}

\section{Lab komputasi: lima pemeriksaan}
\tingkat{lanjut}

Lima klaim di bab ini dapat diperiksa dengan satu skrip: laju $1/\sqrt n$ pada
hukum bilangan besar, kemencengan dan kurtosis rata-rata eksponensial, jarak
supremum ke normal, galat baku Monte Carlo, dan kegagalan total pada Cauchy.
Semua angka di bawah ini adalah keluaran nyata dari eksekusi dengan NumPy 2.5.2
dan SciPy 1.18.0, pembangkit \texttt{default\_rng(20261024)}.

\begin{lab}
\begin{verbatim}
import numpy as np
from scipy import stats

rng = np.random.default_rng(20261024)

print("== 1. hukum bilangan besar: rata-rata dadu ==")
for k in range(1, 6):
    n = 10**k
    m = rng.integers(1, 7, size=(2000, n)).mean(axis=1)
    print("n=%6d  rata2=%.4f  rms=%.4f  sd/akar n=%.4f"
          % (n, m.mean(), np.sqrt(np.mean((m-3.5)**2)),
             np.sqrt(35/12/n)))

print("== 2. CLT: rata-rata eksponensial ==")
for n in (2, 5, 10, 30, 100, 500):
    z = (rng.exponential(size=(50000, n)).mean(axis=1)
         - 1)*np.sqrt(n)
    ks = stats.kstest(z, "norm")
    print("n=%4d skew=%+.3f (%+.3f) kurt=%+.3f (%+.3f)"
          % (n, stats.skew(z), 2/np.sqrt(n),
             stats.kurtosis(z), 6.0/n))
    print("      KS D=%.4f p=%.3f" % (ks.statistic, ks.pvalue))

print("== 3. jarak sup ke normal, eksak lewat gamma ==")
rho = 2 + 2*(6/np.e - 2)          # E|X-1|^3 untuk X ~ Exp(1)
for n in (2, 5, 10, 30, 100, 500):
    zg = np.linspace(-4, 4, 40001)
    Fn = stats.gamma.cdf(n*(1+zg/np.sqrt(n)), a=n)
    D = np.max(np.abs(Fn - stats.norm.cdf(zg)))
    print("n=%4d sup|Fn-Phi|=%.5f  D*akar n=%.3f  BE=%.3f"
          % (n, D, D*np.sqrt(n), 0.4748*rho/np.sqrt(n)))

print("== 4. galat baku Monte Carlo untuk E[exp(U)] ==")
for k in (2, 3, 4, 5):
    n = 10**k
    est = np.exp(rng.random(size=(2000, n))).mean(axis=1)
    print("n=%6d  sd=%.5f  sd*akar n=%.5f"
          % (n, est.std(), est.std()*np.sqrt(n)))

print("== 5. Cauchy: rata-rata tidak memusat ==")
for k in (1, 2, 3, 4):
    n = 10**k
    m = rng.standard_cauchy(size=(4000, n)).mean(axis=1)
    print("n=%5d  median|m|=%.3f P(|m|>10)=%.4f sd=%.2f"
          % (n, np.median(np.abs(m)), np.mean(np.abs(m) > 10),
             m.std()))
print("teori P(|rata2|>10) = %.4f"
      % (1 - 2/np.pi*np.arctan(10)))
\end{verbatim}
\end{lab}

Keluarannya, apa adanya:

\begin{verbatim}
== 1. hukum bilangan besar: rata-rata dadu ==
n=    10  rata2=3.5014  rms=0.5440  sd/akar n=0.5401
n=   100  rata2=3.4992  rms=0.1698  sd/akar n=0.1708
n=  1000  rata2=3.4986  rms=0.0527  sd/akar n=0.0540
n= 10000  rata2=3.5001  rms=0.0171  sd/akar n=0.0171
n=100000  rata2=3.4997  rms=0.0053  sd/akar n=0.0054
== 2. CLT: rata-rata eksponensial ==
n=   2 skew=+1.440 (+1.414) kurt=+3.130 (+3.000)
      KS D=0.0929 p=0.000
n=   5 skew=+0.919 (+0.894) kurt=+1.393 (+1.200)
      KS D=0.0602 p=0.000
n=  10 skew=+0.631 (+0.632) kurt=+0.589 (+0.600)
      KS D=0.0480 p=0.000
n=  30 skew=+0.355 (+0.365) kurt=+0.181 (+0.200)
      KS D=0.0219 p=0.000
n= 100 skew=+0.196 (+0.200) kurt=+0.038 (+0.060)
      KS D=0.0167 p=0.000
n= 500 skew=+0.112 (+0.089) kurt=-0.003 (+0.012)
      KS D=0.0062 p=0.043
== 3. jarak sup ke normal, eksak lewat gamma ==
n=   2 sup|Fn-Phi|=0.09446  D*akar n=0.134  BE=0.811
n=   5 sup|Fn-Phi|=0.05963  D*akar n=0.133  BE=0.513
n=  10 sup|Fn-Phi|=0.04211  D*akar n=0.133  BE=0.363
n=  30 sup|Fn-Phi|=0.02429  D*akar n=0.133  BE=0.209
n= 100 sup|Fn-Phi|=0.01330  D*akar n=0.133  BE=0.115
n= 500 sup|Fn-Phi|=0.00595  D*akar n=0.133  BE=0.051
== 4. galat baku Monte Carlo untuk E[exp(U)] ==
n=   100  sd=0.04928  sd*akar n=0.49281
n=  1000  sd=0.01586  sd*akar n=0.50156
n= 10000  sd=0.00495  sd*akar n=0.49538
n=100000  sd=0.00155  sd*akar n=0.49056
== 5. Cauchy: rata-rata tidak memusat ==
n=   10  median|m|=1.004 P(|m|>10)=0.0630 sd=33.04
n=  100  median|m|=1.008 P(|m|>10)=0.0680 sd=863.61
n= 1000  median|m|=1.040 P(|m|>10)=0.0653 sd=23.87
n=10000  median|m|=0.973 P(|m|>10)=0.0580 sd=379.53
teori P(|rata2|>10) = 0.0635
\end{verbatim}

Bagian 1 memperlihatkan hukum bilangan besar bekerja, dan memperlihatkan
harganya. Kolom rms cocok dengan $\sigma/\sqrt n$ sampai dua angka di belakang
koma untuk $n\ge100$: $0{,}1698$ lawan $0{,}1708$, $0{,}0171$ lawan $0{,}0171$,
$0{,}0053$ lawan $0{,}0054$. Tetapi lihat besarnya: pada $n=10^5$ rata-ratanya
masih meleset sekitar $0{,}005$, yaitu $0{,}15$ persen dari $3{,}5$. Rata-rata
$10^5$ lemparan bukan "praktis tepat"; ia tepat sampai dua angka di belakang
koma.

Bagian 2 dan 3 memperlihatkan dua cara mengukur kedekatan ke normal. Kemencengan
dan kurtosis mengikuti ramalan teori dengan rapi mulai $n=10$: $0{,}631$ lawan
$2/\sqrt{10}=0{,}632$ dan $0{,}589$ lawan $6/10=0{,}600$. Pada $n=2$ dan $n=5$
penaksir kurtosis masih berisik, karena simpangannya bergantung pada momen
kedelapan distribusi itu, bukan hanya pada banyaknya ulangan. Bagian 3 menghitung
jarak supremum secara eksak, tanpa simulasi, dengan memakai sebaran Gamma untuk
rata-rata eksponensial: $0{,}04211$ pada $n=10$ dan $0{,}00595$ pada $n=500$.
Kolom $D\sqrt n$ hampir tetap di $0{,}133$, bukti langsung bahwa lajunya
$1/\sqrt n$. Batas Berry--Esseen $1{,}1464/\sqrt n$ berlaku, dengan angka
$0{,}209$ lawan $0{,}02429$ pada $n=30$; batas itu sekitar $8{,}6$ kali lebih
longgar daripada kenyataan, dan itu wajar untuk batas yang seragam atas semua
distribusi.

Uji Kolmogorov--Smirnov pada lima puluh ribu tarikan menolak kenormalan dengan
$p=0{,}000$ untuk semua $n$ sampai $100$, lalu memberi $p=0{,}043$ pada $n=500$.
Itu bukan tanda program salah: aproksimasi normal tidak pernah tepat untuk $n$
berhingga, dan jaraknya hanya mengecil sampai uji sebesar itu tidak mampu
melihatnya lagi. Pertanyaan yang benar bukan "apakah normal" melainkan
"seberapa besar galat yang masih boleh kita terima", dan bagian 3 menjawabnya
dengan angka.

Bagian 4 memisahkan dua besaran yang sering tertukar. Simpangan baku penaksir
turun dari $0{,}04928$ ke $0{,}00155$ ketika $n$ naik seribu kali, sedangkan
$\text{sd}\cdot\sqrt n$ hampir tetap di $0{,}49$. Nilai tetap itu adalah
$\sigma$ dari $\exp(U)$, yaitu $\sqrt{(e^2-1)/2-(e-1)^2}=0{,}491971$. Inilah
hukum $1/\sqrt n$ dalam bentuk paling telanjang: untuk ketelitian $0{,}001$ pada
$95$ persen keyakinan kita perlu $(1{,}96\cdot0{,}491971/0{,}001)^2=929803$
tarikan.

Bagian 5 adalah tandingannya. Median $|\bar X_n|$ tidak mengecil; peluang
melewati sepuluh tidak mengecil; dan simpangan baku empirisnya melompat dari
$33{,}04$ ke $863{,}61$ lalu turun ke $23{,}87$. Teori menyebut $0{,}0635$, dan
kelima angka simulasi berkisar di situ. Tidak ada hukum bilangan besar di sini,
dan tidak ada teorema limit pusat. Satu-satunya cara menyelamatkan analisis adalah
mengubah pertanyaannya: menaksir indeks ekor, bukan rata-rata.

\begin{untukdosen}
Urutan mengajar yang saya sarankan: mulai dari tabel simulasi, bukan dari definisi.
Tunjukkan pada kelas bahwa galat satu jalankan tidak mengikuti pola apa pun,
sedangkan sebaran galatnya mengikuti. Sesudah itu tiga modus konvergensi, dengan
contoh penyangkal untuk setiap arah yang gagal; jangan menunda contoh-contoh itu,
karena di situlah mahasiswa paham bahwa "menuju" perlu diperjelas. Baru kemudian
Markov, Chebyshev, hukum lemah, hukum kuat sebagai pernyataan, lalu limit pusat
dengan bukti fungsi pembangkit momen.

Dua miskonsepsi yang paling sering muncul. Pertama, mahasiswa membaca "konvergen
dalam peluang" sebagai "nilainya menjadi tetap", lalu heran ketika simulasi
menunjukkan lompatan besar pada $n$ besar; contoh $(0,1)$ yang dipecah pada bab
ini adalah penawarnya. Kedua, mereka membaca teorema limit pusat sebagai "rata-rata
menjadi normal", padahal yang menjadi normal adalah galat yang dibagi
$\sigma/\sqrt n$; akibatnya mereka menuliskan $\bar X_n\approx N(\mu,\sigma^2)$
dan lupa $\sqrt n$.

Aktivitas papan tulis yang berhasil: minta dua mahasiswa memerankan penjudi dan
bandar. Penjudi menyebut toleransi, bandar menyebut banyak lemparan yang
menjaminnya, dan kelas menghitung dengan Chebyshev lalu membandingkannya dengan
angka eksak dari SciPy. Selisihnya selalu besar, dan itulah pintu masuk ke
aproksimasi normal. Tanda kelas siap lanjut: mereka dapat menjelaskan mengapa
Chebyshev tetap berguna walaupun longgar, dan dapat menyebut satu distribusi yang
membuat teorema limit pusat gagal.
\end{untukdosen}

\section{Ringkasan}

\begin{ringkasan}
\begin{itemize}
  \item Ada tiga modus konvergensi yang berbeda: hampir pasti, dalam peluang, dan
        dalam distribusi. Hampir pasti mengakibatkan dalam peluang, dalam peluang
        mengakibatkan dalam distribusi, dan arah sebaliknya gagal kecuali bila
        limitnya konstanta.
  \item Ketaksamaan Markov dan Chebyshev hanya memakai momen, karena itu berlaku
        untuk semua distribusi bervariansi hingga, dan karena itu pula longgar.
  \item Hukum bilangan besar lemah: $P(|\bar X_n-\mu|\ge\varepsilon)\le
        \sigma^2/(n\varepsilon^2)\to0$. Hukum kuat menambahkan bahwa
        $\bar X_n\to\mu$ hampir pasti bila $\mathbb{E}|X_1|<\infty$, tanpa
        menuntut variansi.
  \item Teorema limit pusat: $\sqrt n(\bar X_n-\mu)/\sigma\rightsquigarrow N(0,1)$
        bila variansi hingga. Lajunya $1/\sqrt n$, maka empat kali data untuk
        separuh galat, dan seratus kali data untuk satu digit tambahan.
  \item Aproksimasi normal untuk jumlah diskret memerlukan koreksi kontinuitas
        setengah satuan; koreksi itu memangkas galat di ekor sekitar dua puluh
        kali pada contoh binomial di bab ini.
  \item Bila variansi tak hingga, teorema limit pusat gagal, bukan hanya melambat.
        Rata-rata $n$ pengamatan Cauchy tetap Cauchy, dan yang menggantikan
        lonceng adalah sebaran stabil berindeks $\alpha<2$.
\end{itemize}
\end{ringkasan}

\section{Latihan}

\begin{latihan}
\begin{enumerate}
  \item Buktikan: bila $X_n\to X$ dan $Y_n\to Y$ dalam peluang, maka
        $X_n+Y_n\to X+Y$ dalam peluang. Tuliskan pertidaksamaan himpunan yang
        dipakai, lalu jelaskan mengapa $\varepsilon/2$ muncul.
  \item Buktikan bahwa ketaksamaan Chebyshev tajam: untuk $0<\sigma^2\le a^2$,
        ambil $X$ dengan $P(X=a)=P(X=-a)=\sigma^2/(2a^2)$ dan
        $P(X=0)=1-\sigma^2/a^2$, lalu hitung $\operatorname{Var}(X)$ dan
        $P(|X|\ge a)$.
  \item Nilai klaim ini: "Karena teorema limit pusat berlaku untuk semua
        distribusi, rata-rata tiga puluh pengamatan selalu boleh dianggap
        normal." Benar, salah, atau belum cukup bukti? Sebutkan syarat yang
        dilanggar atau yang belum diperiksa, dan pakai angka dari bagian lab.
  \item Untuk koin dengan $p=1/2$ dan toleransi $\varepsilon=0{,}05$, cari $n$
        terkecil yang dijamin Chebyshev untuk peluang meleset paling besar
        $0{,}05$. Lalu cari $n$ terkecil yang benar-benar memenuhi syarat itu,
        dengan menghitung sebaran binomial di Python. Bandingkan kedua angka.
  \item Buktikan bahwa konvergensi dalam peluang tidak menjamin konvergensi
        harapan, dengan memakai $P(X_n=n^2)=1/n$ dan $P(X_n=0)=1-1/n$. Hitung
        $\mathbb{E}[X_n]$ dan $P(|X_n|>\varepsilon)$ untuk $\varepsilon$ tetap.
  \item Untuk $S\sim\text{Binomial}(200;\,0{,}3)$, aproksimasi $P(S\le50)$ dengan
        normal tanpa dan dengan koreksi kontinuitas, lalu bandingkan keduanya
        dengan nilai eksak dari SciPy. Laporkan kedua galat.
  \item Bangkitkan $4000$ rata-rata dari $n$ bilangan Cauchy untuk
        $n=10,10^2,10^3,10^4$, lalu laporkan median $|\bar X_n|$, peluang
        $|\bar X_n|>10$, dan simpangan bakunya. Bandingkan dengan
        $1-(2/\pi)\arctan10$ dan jelaskan mengapa kolom simpangan baku tidak
        mengecil.
  \item Buktikan dengan fungsi karakteristik bahwa rata-rata $n$ variabel iid
        Cauchy berdistribusi Cauchy. Lalu jelaskan mengapa laporan berbentuk
        "rata-rata $\pm$ dua galat baku" pada data semacam itu tidak dapat
        ditafsirkan.
  \item Berapa tarikan Monte Carlo yang diperlukan agar penaksir
        $\mathbb{E}[\exp(U)]$ memiliki setengah lebar interval $95$ persen
        sebesar $0{,}0005$? Pakai $\sigma=0{,}491971$ dari bagian lab, lalu
        bandingkan dengan $n$ yang diperlukan untuk setengah lebar $0{,}001$.
        Nyatakan rasionya dan jelaskan asalnya.
  \item Buktikan bahwa bila $X_n\rightsquigarrow c$ dengan $c$ konstanta, maka
        $X_n\to c$ dalam peluang, tanpa memakai bagian (c) teorema hubungan tiga
        modus di awal bab, lalu tunjukkan dengan contoh bahwa syarat "konstanta"
        tidak dapat dibuang.
\end{enumerate}

Petunjuk untuk tiga butir tersulit. Butir 1: tulis
$\{|(X_n+Y_n)-(X+Y)|>\varepsilon\}$ termuat dalam gabungan dua kejadian, lalu
pakai batas gabungan. Butir 8: fungsi karakteristik jumlah adalah hasilkali
fungsi karakteristik, dan fungsi karakteristik Cauchy $e^{-|t|}$ bersifat
istimewa karena bentuknya tidak berubah saat dipangkatkan. Butir 10: tulis
$P(X_n\le t)\le P(X_n\le c-\varepsilon)+P(|X_n-c|>\varepsilon)$ untuk $t$ di
kiri $c$, dan ingat bahwa sebaran titik massa melompat tepat di $c$.
\end{latihan}

\begin{untukpeneliti}
Yang benar-benar terbukti di bab ini: implikasi hampir pasti $\Rightarrow$ dalam
peluang $\Rightarrow$ dalam distribusi, hukum lemah lewat Chebyshev, dan limit
pusat lewat fungsi pembangkit momen dengan asumsi $\mathbb{E}|Y_1|^3<\infty$.
Yang populer tetapi tidak terbukti di sini: hukum kuat (pernyataan saja, bukti
Etemadi dilewati), versi limit pusat tanpa momen ketiga (butuh fungsi
karakteristik dan teorema kontinuitas Levy), dan seluruh teori sebaran stabil
(pernyataan Gnedenko--Kolmogorov, tanpa bukti).

Sumber primer yang saya pakai: Wasserman, \emph{All of Statistics}, bab 5 bagian
5.2 (tipe konvergensi dan diagram implikasinya), 5.3 (hukum lemah), 5.4 (limit
pusat dan Berry--Esseen), serta lampiran 5.7.1 dan 5.7.2 untuk hukum kuat dan
bukti limit pusat. Blitzstein dan Hwang, \emph{Introduction to Probability}, bab
10 bagian 10.1 (ketaksamaan), 10.2 (hukum bilangan besar), 10.3 (limit pusat).
Untuk sisi ekor berat: Feller jilid 2.

Batas keberlakuan yang perlu disebut. Pertama, semua teorema di sini mengandaikan
kebebasan dan distribusi sama; untuk data dependen, misalnya deret waktu atau
iterasi rantai Markov, yang berlaku adalah versi martingale atau mixing, dan
lajunya bisa jauh lebih lambat. Kedua, konstanta optimal pada Berry--Esseen belum
diketahui; $0{,}4748$ cukup, dan batas bawah yang terbaik bagi konstanta itu
masih di bawah $0{,}5$. Ketiga, laju $1/\sqrt n$ adalah laju paling buruk yang
dijamin secara seragam; untuk distribusi tertentu laju sebenarnya lebih baik,
dan untuk ekor berat dengan $\alpha\in(1,2)$ laju sebenarnya lebih buruk.
Pertanyaan terbuka yang masih hidup: seberapa tajam batas Berry--Esseen untuk
kelas distribusi yang dibatasi momennya, dan bagaimana menaksir indeks ekor
dengan interval kepercayaan yang sah pada $n$ kecil.
\end{untukpeneliti}

\peta{Bab 19 pada Wasserman, \emph{All of Statistics}: bab 5 (Convergence of
Random Variables), bagian 5.2 Types of Convergence, 5.3 The Law of Large Numbers,
5.4 The Central Limit Theorem, dan lampiran 5.7.1 Almost Sure and L1 Convergence
serta 5.7.2 Proof of the Central Limit Theorem. Pada Blitzstein dan Hwang,
\emph{Introduction to Probability}: bab 10 (Inequalities and limit theorems),
bagian 10.1 Inequalities, 10.2 Law of large numbers, 10.3 Central limit theorem.
Untuk sebaran stabil: Feller, \emph{An Introduction to Probability Theory and Its
Applications}, jilid 2. Distribusi dasar yang dipakai di sini dibangun di Bab 17,
dan interval kepercayaan yang bersandar pada bab ini dibahas di Bab 19.}
